The products derived from the bowling release are the ball’s 2D trajectory along the lane and its spin. The trajectory is defined by localizing the geometric center of the ball through the clip’s frames, whereas the parameters of spin are axis orientation and angular velocity. The previously obtained data is used to define the region of interest in which we’ll find the ball’s projection as a circle in every frame. Once the trajectory data is available, the ball’s locations across consecutive frames are used to determine the positional change of its features and to derive its spin parameters. With both products in hand, we can then generate a visualization that consistently follows the ball’s motion and proceed with the 3D reconstruction of the scene.

\subsubsection{Trajectory detection}

The ball-tracking workflow consists of three main steps: video preprocessing, ball-candidate detection, and trajectory computation. We begin by isolating the scene’s region of interest and applying noise-reduction and contrast-enhancement filters to each frame. Then, the \emph{Circle Hough Transform} (CHT) algorithm is applied to generate a dataset of all geometries that could represent the ball’s position in each frame. Finally, we select the geometries that best form a coherent trajectory from the candidates, interpolate any missing values in the set, and smooth the final trajectory. In the end, we obtain a ball-center dataset with the image geometric coordinates and their corresponding radii.

Firstly, each frame must be manipulated to reduce video noise so that the CHT consistently detects the ball. This process starts by masking all pixels outside the ROI using the coordinates retrieved in section \ref{sec:lane_detection}. Then, median blurring is applied to remove high-frequency noise from light reflections in the lane, followed by adaptive histogram equalization to increase contrast specifically in the region where the ball is located. This last step is particularly helpful because the ROI is expected to have a relatively consistent surface across the frames. We end up with a black-and-white clip in which the ball is clearly distinguishable from the rest of the scene, and the lane reflections are smoothed out.

\begin{figure}[H]
    \centering
    \includegraphics[height=0.27\textwidth]{ball_tracking/preprocessed_frame_clip_1.png}
    \includegraphics[height=0.27\textwidth]{ball_tracking/preprocessed_frame_clip_2.png}
    \includegraphics[height=0.27\textwidth]{ball_tracking/preprocessed_frame_clip_13.png}
    \caption{Preprocessed frames of different bowling releases. The ball is clearly distinguishable from the rest of the scene, and the lane imperfections are smoothed out.}
    \label{fig:preprocessed_frame}
\end{figure}


Once we have the preprocessed clip, we apply OpenCV's built-in CHT algorithm. The library implements the \emph{Hough gradient method}, which first applies \emph{Canny edge detection} and then computes the local gradient for every nonzero pixel via the first-order Sobel x and y derivatives. These gradients are used to populate the two-dimensional \emph{accumulator plane} \cite{BRADSKI08}, in which potential circle centers are tracked. For every edge pixel, a line is projected along its gradient from a specified minimum to a maximum distance, incrementing the accumulator cells along the way. The local maxima in this plane are selected as candidate centers. Finally, for each center, the algorithm evaluates the distances to nearby edge pixels to determine the optimal radius, preserving the circle only if it receives sufficient support and maintains a minimum distance from previously validated centers. The result is a dataset containing every detected circle, its x- and y-center coordinates, and its radius for each frame.

\begin{figure}[H]
    \centering
    \includegraphics[height=0.2\textwidth]{ball_tracking/candidates_frame_1.png}
    \includegraphics[height=0.2\textwidth]{ball_tracking/candidates_frame_2.png}
    \includegraphics[height=0.2\textwidth]{ball_tracking/candidates_frame_4.png}
    \includegraphics[height=0.2\textwidth]{ball_tracking/candidates_frame_3.png}
    \caption{Detected circles overlayed on the preprocessed frames of the same clip. The algorithm detects multiple circles in each frame, which may include false positives due to noise.}
    \label{fig:candidates_frame}
\end{figure}

This algorithm detects circles across video frames; however, it may also introduce noise due to numerical instabilities \cite{BRADSKI08}, leading to false detections, or even no detections over the ball, as seen in figure \ref{fig:candidates_frame}. The third step of the workflow focuses on filtering all circles representing the ball in each frame, resulting in a final defined trajectory. The methodology for discernment was inspired by \cite{PALLAVI08}, in which they elaborate a weighted directed graph to determine the ball’s path in a soccer broadcast. The difference in our implementation is that a bowling release is less dynamic than a soccer kick; its behavior is more predictable, and its trajectory is easier to generalize.

The main idea is that each ball candidate in a frame is represented by a graph node, which is then directionally connected to all other candidates in the five subsequent frames. This is done because the previous step might miss the ball in consecutive frames. In the end, the longest path in the graph is retrieved, which represents the actual ball trajectory. To fill in the missing frames in the trajectory, cubic spline interpolation is applied to both the x- and y-center coordinates, followed by a light Gaussian filter ($\sigma = 5$) to eliminate noise, and a simple quadratic least squares interpolation for the radius values.

\begin{figure}[H]
    \centering
    \includegraphics[width=1\textwidth]{ball_tracking/graph.png}
    \caption{Graph representation of the ball candidates across frames. Each node represents a candidate circle, and directed edges connect candidates in consecutive frames. The longest path through the graph corresponds to the most coherent trajectory of the ball.}
    \label{fig:graph}
\end{figure}

The entire workflow yields the final dataset, which contains the coordinates and radius of each circle from the first to the last detected frame of the clip. This data will be crucial for the following steps in spin detection and 3D reconstruction.

\begin{figure}[H]
    \centering
    \includegraphics[height=0.4\textwidth]{ball_tracking/comparative_trajectory.png}
    \caption{Left: Observed ball centers (blue) and their corresponding radii (orange) across frames. Right: Final trajectory after interpolation and smoothing, with the same color scheme.}
    \label{fig:comparative_trajectory}
\end{figure}


\subsubsection{Spin detection}

Once we have defined the bounding box, we select the trackable features on the ball’s surface. To compute optical flow, we identify points in the first image that are sufficiently distinct so they can be located again in the next image. OpenCV provides a pipeline based on the Shi-Tomasi corner detection algorithm \cite{SHI94}. Instead of relying solely on edge information, which captures gradient changes in a single direction, the algorithm seeks regions with significant intensity variation in two orthogonal directions. It computes the first-order image gradients, $I_x$ and $I_y$, using a 3x3 Sobel kernel, which represent the rate of change of brightness in the horizontal and vertical directions. These gradients are then used to compute the \emph{structure tensor} (or second-moment matrix) $M$ over a window function $w$ that captures the spatial context around a pixel:

\begin{equation}
	M = \sum_{x,y}w(x, y) 
\begin{bmatrix} 
	I^2_x & I_x I_y \\
	I_x I_y & I^2_y
\end{bmatrix}
\end{equation}

The structure tensor is a covariance matrix of the image gradients. Its eigenvalues $\lambda_1$ and $\lambda_2$ are the intensity change in each direction. If both are small, the region is flat; if one is larger than the other, it represents an edge; and if both are large, we have found the so-called corner, where there is a strong intensity change in both orthogonal directions. Corners typically represent the best trackable features for the optical flow. In this context, the intensity changes must be detected solely on the ball’s surface. Hence, the importance of the ball not having a homogeneous texture. We finally end up with an array of pixel locations that hopefully will be found in the next frame.

With the initial features detected, we approach their displacement between consecutive frames using the \emph{Pyramidal Lucas-Kanade} (PyLK) optical flow algorithm \cite{BOUGUET00}. The standard Lucas-Kanade algorithm rests on three core assumptions about the scene: brightness constancy, spatial coherence, and temporal persistence. The first implies that the brightness of tracked pixels does not change between frames. Spatial coherence means that neighboring points in a scene belong to the same surface and have similar motion. The last assumption is that features move relatively small distances between consecutive frames, as it relies only on a local window of information \cite{BRADSKI08}. The latter, however, conflicts with the behavior of the spinning ball: the ball’s rotational speed causes features to move noticeably between frames.

We want to determine the velocity vectors that indicate how far the feature has moved in the $x$ and $y$ directions. The brightness constancy tells us that the pixel intensity $I(x, y, t)$ at a specific point on the bowling ball does not change as the ball moves to a new position in the next frame:

\begin{equation}
I(x, y, t) = I(x + \Delta x, y + \Delta y, t + 1)
\end{equation}

By linearizing, we obtain the \emph{standard optical flow equation}:

\begin{equation}
I_xu + I_yv = -I_t
\end{equation}

, where $I_t$ is the difference in pixel intensity between frames, $I_x$ and $I_y$ are the spatial image gradients (computed in the feature detection), and $u$ and $v$ are the unknown velocity vectors. However, as we have a single equation for two unknowns, it cannot be solved with a single pixel. The solution takes advantage of the spatial coherence assumption: the optical flow is calculated for all the pixels around the feature for which we can build a system solvable by least-squares minimization:

\begin{equation}
    \overset{A}{\
    \begin{bmatrix}
        I_x(p_1) & I_y(p_1) \\
        I_x(p_2) & I_y(p_2) \\
        \vdots & \vdots \\
        I_x(p_n) & I_y(p_n)
    \end{bmatrix}}_{[n\times 2]}
    \overset{x}{\
    \begin{bmatrix}
        u \\
        v \\
    \end{bmatrix}}_{[2\times 1]}=
    \overset{b}{\
    \begin{bmatrix}
        I_t(p_1) \\
        I_t(p_2) \\
        \vdots \\
        I_t(p_n)
    \end{bmatrix}}_{[n\times 1]}
\end{equation}

